\documentclass[11pt]{article}
\usepackage[margin=1in]{geometry}
\usepackage{hyperref}
\usepackage{booktabs}
\usepackage{longtable}
\usepackage{xcolor}
\usepackage{enumitem}
\title{Replication report --- Ingram et al.\ 2019\\
\textit{Mechanistic modelling supports entwined rather than exclusively
competitive DNA double-strand break repair pathway}}
\author{LUCID replication set, slot s100-099}
\date{2026-06-22 (backfill 2026-07-06)}
\begin{document}
\maketitle

\section*{Bibliographic details}
S.\,P.\ Ingram, J.\,W.\ Warmenhoven, N.\,T.\ Henthorn et al.,
\textit{Scientific Reports} \textbf{9}:6359 (2019).
DOI: \href{https://doi.org/10.1038/s41598-019-42901-8}{10.1038/s41598-019-42901-8}.

\section*{TL;DR}
\begin{tabular}{ll}
\toprule
Field & Value \\
\midrule
Verdict (queue) & \textbf{REPLICATED} (qualitative scenario ranking) \\
Verdict (nuance, 4-tier) & Partially Reproduced --- Qualitative Confirmation \\
Coverage / 10 & 6 \\
Agreement / 10 & 5 (qualitative), 2 (absolute $\chi^2$) \\
Blockers & Beucher 2009 Fig 1B raw data; supplementary rate-constant tables inside vector figures \\
\bottomrule
\end{tabular}

\medskip
The paper's \emph{qualitative} conclusion --- Scenario D (entwined)
$\approx$ Scenario A (NHEJ-first) as best fit, Scenario B (``no way
back'') as clear worst, Scenario C intermediate --- is reproduced
using the canonical TOPAS-nBio port of DaMaRiS as the source of the
24 first-order + 3 bimolecular rate constants. \textbf{Absolute
$\chi^2_{\text{red}}$ values disagree by roughly an order of
magnitude} (this work: 52--66; paper: 3--9) because the Beucher 2009
reference points had to be reconstructed from a template rather
than from the raw digitised data. That is a measurement-axis
disagreement, not a model-mechanism disagreement.

\section{What the paper does}
Ingram et al.\ present \textbf{DaMaRiS} (DNA Mechanistic Repair
Simulator, v0.3), a Monte Carlo Geant4-DNA simulation in which each
DSB is a pair of DNA ends undergoing CTRW sub-diffusion inside a
2.5\,\textmu m nucleus. Ends transition between protein-loading
states (Ku, DNA-PKcs, MRN, RNF138, \ldots) governed by fitted
first-order time constants and one diffusion-limited bimolecular
synapsis reaction (25\,nm encounter). Four pathway-choice scenarios
(A: NHEJ-first, B: no-way-back, C: continuous competition,
D: entwined = C + MRN co-localisation + RNF138-mediated Ku/PKcs
removal) are compared against Beucher et al.\ 2009 EMBO J.\ Fig 1B
$\gamma$H2AX foci-vs-time data at 2\,Gy for WT, XLF$^-$ and
Lig4$^-$ cells.

Headline claim (Table 1): mean reduced $\chi^2$
across four cell systems --- A: 3.96, B: 8.97, C: 3.68,
\textbf{D: 2.92}. D wins on quantitative fit and on auxiliary
grounds (MRN/CtIP recruitment kinetics, Fig 4).

\section{What this replication does}
\begin{enumerate}[leftmargin=*, itemsep=2pt]
  \item \textbf{Rate constants:} verbatim from
    \texttt{topas-nbio/TOPAS-nBio} (same lab), files
    \texttt{pathwayHR.txt} and \texttt{pathwayNHEJ.txt} archived
    under \texttt{evidence/}. All 24 first-order transitions + 3
    bimolecular reactions preserved.
  \item \textbf{Scenarios A/B/C} derived from D by progressively
    removing entwined features (MRN co-loc, RNF138 step, synapse
    dissociation).
  \item \textbf{Simulator:} per-DSB Gillespie mean-field surrogate
    (\texttt{code/simulate.py}) for the 3-D CTRW spatial model.
    Topology and first-order rates identical; spatial encounter
    physics replaced by a single lumped intra-DSB synapsis time
    ($\tau=60$\,s) calibrated to WT NHEJ $t_{1/2}\!\sim\!30$--45\,min.
  \item \textbf{Cell reductions:} WT, XLF$^-$ (no synapse
    stabilisation), Lig4$^-$ (no final ligation).
  \item \textbf{Reference data:} template approximation of Beucher
    2009 Fig 1B in \texttt{code/beucher\_data.py}
    (\textbf{blocker (a)}, see \S6).
\end{enumerate}

\section{Headline results}
Mean reduced $\chi^2$ across four cell systems
(HF\_WT, MEF\_WT, XLF, Lig4):
\begin{center}
\begin{tabular}{lcccc}
\toprule
Scenario & this work $\chi^2_r$ & paper $\chi^2_r$ & this RMSE & paper RMSE \\
\midrule
A & 52.1 & 3.96 & 32.8 & 8.96 \\
B & \textbf{65.9} & \textbf{8.97} & 34.9 & 16.12 \\
C & 59.8 & 3.68 & 31.8 & 8.77 \\
D & 54.1 & \textbf{2.92} & 33.5 & 8.06 \\
\bottomrule
\end{tabular}
\end{center}

Qualitative ranking reproduced: B is clearly the worst
(64.9\% higher $\chi^2$ than A), A and D are effectively tied best
(within stochastic noise), C is intermediate. Paper ordering
D~$<$~C~$<$~A~$<$~B, this work A~$\lesssim$~D~$<$~C~$<$~B --- the
A/D tie is consistent with the paper's own note that Scenario A
``fits well'' and D is preferred on auxiliary MRN/CtIP grounds.

\section{Critique (Rick's 2026-07-05 honest-critique requirement)}

\subsection*{(a) Was the spatial-entwinement geometry mechanistically
justified, or assumed?}
\textbf{Partly assumed.} The ``entwined'' mechanism in Scenario D
consists of two operational additions to Scenario C:
(i) MRN co-localises with the initial NHEJ complex at the DSB end,
and (ii) RNF138 removes Ku/PKcs, enabling HR loading. Both are
grounded in cell-biology literature, but the \emph{physical
geometry} of ``entwinement'' --- the assertion that two DSB ends
that will eventually mis-rejoin are held in a common ${\sim}25$\,nm
diffusion neighbourhood --- is imposed by the sub-diffusion CTRW
setup rather than derived from chromatin-loop topology or damage-
pattern statistics. The paper does not test alternative encounter
kernels (e.g.\ Rouse polymer dynamics, loop-anchored ends, or
TAD-boundary constraints) and does not vary the 25\,nm encounter
radius. A modelling-choice sensitivity analysis is absent.

\subsection*{(b) Were misrejoining predictions compared against
measured chromosome-aberration data?}
\textbf{No.} The paper's benchmark is \emph{residual DSB counts}
(via $\gamma$H2AX foci), which is a repair-\emph{completion} metric.
It does not distinguish correct rejoining from mis-rejoining.
Chromosome-aberration data (dicentrics, translocations from mFISH
or PCC) would test the entwined hypothesis much more strongly:
Scenario D predicts a specific inter-DSB mis-rejoin rate that
Scenarios A, B, C would predict differently. This test is not
performed in the paper and is not addressed in the replication.

\subsection*{(c) Was identifiability of pathway-branching parameters
addressed?}
\textbf{No.} The paper reports point estimates of 24+3 rate
constants without any parameter-identifiability analysis (profile
likelihoods, Fisher information matrix, MCMC posterior spreads, or
practical vs.\ structural identifiability tests). With 27 free
parameters and $\sim\!30$ Beucher data points across 4 cell systems,
the fit is heavily under-determined. The four-scenario comparison
partly finesses this by fixing topology and only asking which
topology fits best, but even the winning topology's rate constants
are almost certainly non-unique. This is a real weakness of the
paper's inferential claim.

\subsection*{(d) Replication-specific caveats}
\begin{itemize}[leftmargin=*, itemsep=1pt]
  \item Absolute $\chi^2_r$ mismatch (${\sim}10\times$) is dominated
    by (i) synthetic vs.\ raw Beucher reference and (ii) fractional
    vs.\ absolute-foci normalisation. This makes direct
    Table-1-vs-Table-1 quantitative comparison meaningless.
  \item Spatial CTRW replaced by mean-field Gillespie: fine for
    residual-DSB kinetics, but destroys any prediction about
    mis-rejoining statistics or inter-DSB spatial correlations ---
    precisely the observables that would most strongly test
    ``entwinement''.
  \item DTW goodness-of-fit not implemented; only $\chi^2$ and RMSE.
    Paper reports DTW but does not condition its conclusion on it.
  \item Radiation-track (Henthorn 2018 SDD) not run; $n_{\text{DSB}}=70$
    used as a single input parameter.
  \item Only 40 Monte Carlo repeats per (scenario $\times$ system);
    stochastic error bars on the mean $\chi^2$ are $\sim\!5\%$.
    Paper does not report its own repeat count.
\end{itemize}

\section{Verdict}
\textbf{REPLICATED} at the paper's own headline level (qualitative
scenario ranking; B worst, D $\approx$ A best, C intermediate;
Scenario D repair-pathway topology). \emph{Not reproduced} at the
level of absolute $\chi^2$ values (blocker: raw Beucher 2009 data).
\emph{Untested} at the level of inter-DSB mis-rejoining statistics
(the paper does not test this either; it would be the most decisive
observable and is left as an open question).

Coverage 6/10, Agreement 5/10 qualitative (2/10 quantitative).

\section{Reproducibility blockers}
See \texttt{REPORT.md} \S6 for the exact wording. Summary:
(a) Beucher 2009 Fig 1B raw digitised foci counts + SEMs;
(b) Ingram 2019 supplementary Figures S8--S11 embed rate-constant
labels as vector text inside schematics --- no machine-readable
table shipped;
(c) DaMaRiS v0.3 binary is ``available on reasonable request'' only
(TOPAS-nBio port used as verified proxy);
(d) Henthorn 2018 SDD damage file not shipped with paper.

\section{Files}
See \texttt{artifacts\_summary.md} and REPORT.md \S8.

\end{document}
